% !TEX program = xelatex
\documentclass{英语试卷}
\begin{document}
\examtitle{诸暨市 2025 届高三第一次诊断性考试试题}{英语}
\par 诸暨市 2025 届高三第一次诊断性考试试题\par
\par 英语学科 试题卷\par
\begin{notice}
\par 1. 答卷前，考生务必将自己的姓名、准考证号填写在答题卡上，\par
\par 2. 回答选择题时，选出每小题答案后，用铅笔把答题卡上对应题目的答案标号涂黑。如需改动，用橡皮擦干净后，再选涂其他答案标号。回答非选择题时，将答案写在答题卡上，写在本试卷上无效。\par
\par 3. 考试结束后，将本试卷和答题卡一并交回。\par
\end{notice}
\parttitle{第一部分 听力（共两节，满分 30 分）}
\par 做题时，先将答案标在试卷上，录音内容结束后，你将有两分钟的时间将试卷上的答案转涂到答题卡上。\par
\sectitle{第一节（共5小题：每小题1.5分，满分7.5分）}
\par 听下面 5 段对话。每段对话后有一个小题，从题中所给的 A、B、C 三个选项中选出最佳选项。听完每段对话后，你都有 10 秒钟的时间来回答有关小题和阅读下一小题。每段对话仅读一遍。\par
\question{1. Where are the speakers?}
\begin{choices}(3)
  \item At home.
  \item In a bookstore.
  \item In a fruit shop.
\end{choices}
\question{2. Why cannot the woman go dancing this evening?}
\begin{choices}(3)
  \item She is dizzy.
  \item She is occupied.
  \item She is exhausted.
\end{choices}
\question{3. When will the magazine probably arrive?}
\begin{choices}(3)
  \item Wednesday.
  \item Thursday.
  \item Friday.
\end{choices}
\question{4. What will the woman do?}
\begin{choices}(3)
  \item Go to the lab with the man.
  \item Work on her experiments.
  \item Check on what's for dinner.
\end{choices}
\question{5. What does the woman mean?}
\begin{choices}(3)
  \item Lucy is a green hand.
  \item She doesn't like the cake.
  \item She has a sweet tooth.
\end{choices}
\sectitle{第二节（共15小题：每小题1.5分，满分22.5分）}
\par 听下面 5 段对话或独白。每段对话或独白后有几个小题，从题中所给的 A、B、C 三个选项中选出最佳选项。听每段对话或独白前，你将有时间阅读各个小题，每小题 5 秒钟；听完后，各小题将给出 5 秒钟的作答时间。每段对话或独白读两遍。\par
\par 听第 6 段材料，回答第 6 至第 7 题。\par
\question{6. What do you know about the motorbike?}
\begin{choices}(3)
  \item It is being repaired.
  \item It hasn't been fixed yet.
  \item It is beyond repair.
\end{choices}
\question{7. What causes the delay?}
\begin{choices}(3)
  \item The spare parts ordered haven't arrived yet.
  \item They forgot to take care of the bike in time.
  \item The brand is unfamiliar to the mechanics.
\end{choices}
\par 听第 7 段材料，回答第 8 至第 10 题。\par
\question{8. What is the man doing now?}
\begin{choices}(3)
  \item Taking a sea trip.
  \item Swimming in the sea.
  \item Having a boating competition
\end{choices}
\question{9. How does the man most probably feel now?}
\begin{choices}(3)
  \item Panicked.
  \item Lonely.
  \item Delighted.
\end{choices}
\question{10. What does the man expect first when getting to Australia?}
\begin{choices}(3)
  \item Relaxation.
  \item Interviews.
  \item Celebrations.
\end{choices}
\par 听第 8 段材料，回答第 11 至第 13 题。\par
\question{11. Why was the man absent from the engineering class?}
\begin{choices}(3)
  \item He knows well the topic.
  \item He was on sick leave.
  \item He had other fish to fry.
\end{choices}
\question{12. What do you know about DME?}
\begin{choices}(3)
  \item It is a conventional fuel.
  \item It is a clean energy.
  \item It's a rare resource.
\end{choices}
\question{13. Why hasn't DME been on the market?}
\begin{choices}(3)
  \item It does damage to the engine.
  \item It's no better than other fuels.
  \item It's pricey to mass-produce.
\end{choices}
\par 听第 9 段材料，回答第 14 至第 17 题。\par
\question{14. When did Ben grow interested in postcards?}
\begin{choices}(3)
  \item When he found a collection at home.
  \item When he met with his grandfather.
  \item When he was told about their value.
\end{choices}
\question{15. Who initially collected postcards in Ben's family?}
\begin{choices}(3)
  \item His mother.
  \item His uncle.
  \item His grandfather.
\end{choices}
\question{16. Which postcards are Ben's favorite?}
\begin{choices}(3)
  \item The old ones with pictures of world-famous cities
  \item Those with photos of famous buildings and harbors.
  \item The ones with pictures of ways people used to travel.
\end{choices}
\question{17. What do both speakers think of the out-of-date postcards?}
\begin{choices}(3)
  \item They are plain-looking.
  \item They are worthless.
  \item They are informative.
\end{choices}
\par 听第 10 段材料，回答第 18 至第 20 题。\par
\question{18. What do you know about Stardate 3000?}
\begin{choices}(3)
  \item The fiction series won't be on tonight.
  \item Viewers have mixed comments on it.
  \item It turns out to be nothing but a failure.
\end{choices}
\question{19. What is Miller and Edwards probably intended to do?}
\begin{choices}(3)
  \item To introduce actors.
  \item To entertain viewers.
  \item To provide directions.
\end{choices}
\question{20. Who is Andy MaKay?}
\begin{choices}(3)
  \item The director of An Insect's Life.
  \item The host of the TV program.
  \item A character in Policewatch.
\end{choices}
\parttitle{第二部分 阅读理解（共两节，满分 50 分）}
\sectitle{第一节（共15小题：每小题2.5分，满分37.5分）}
\par 阅读下列短文，从每题所给的 A、B、C 和 D 四个选项中，选出最佳答案。\par
\passagelabel{A}
\par Best Washington, D. C. Museums\par
\par Washington, D. C. has wonderful museums that appeal to virtually any visitor's interest. From large publicly funded institutions to small historic homes, get ready to learn a great deal and find some welcome surprises and bonus.\par
\par Smithsonian National Museum of Natural History\par
\par There is something for everyone at this world-famous museum. It houses so many artifacts(手工艺品) that you can't possibly see them all in one visit. The dinosaur exhibits are fascinating and great for kids. Then there's amazing scenery in the IMAX films.\par
\par Visiting Tip: If you are visiting with kids, be sure to have a shot at hands-on activities.\par
\par Mount Vernon Estate and Gardens\par
\par Located outside of the city, the home of George Washington is often missed. The estate is the most scenic tourist attraction in the Washington, D. C. area. The historic residence is restored and decorated as it was when Washington lived there. The museum and education center features25 state-of-the-art galleries and theaters that tell the story of Washington's life.\par
\par Visiting Tip: Visit the Mansion first as it is usually the most crowded.\par
\par U.S. Holocaust Memorial Museum\par
\par The museum is a memorial to the millions of Jews who died during the Nazi regime in Germany during World War II. The exhibits tell the horrific story of genocide and teach the dangers of hatred and prejudice. Visiting this museum is an emotional experience. The permanent exhibits are not recommended for children under 11.\par
\par Visiting Tip: Timed passes are distributed for the same day on a first come, first served basis.\par
\par International Spy Museum\par
\par Kids love looking at spy devices and cameras at International Spy Museum, not to mention the interactive spy missions you can try out here.\par
\par Visiting Tip: Tickets for the museum start at 22.95 for adults, 16.95 for seniors, and \$14.95 for children ages 7 to 11, and children under 6 are free.\par
\question{21. Which of the four museums is unsuitable for children under 11?}
\begin{choices}(2)
  \item Mount Vernon Estate and Gardens
  \item U. S. Holocaust Memorial Museum
  \item International Spy Museum
  \item National Museum of Natural History
\end{choices}
\question{22. What can visitors gain by visiting Mount Vernon Estate and Gardens?}
\begin{choices}(2)
  \item A feast of ancient architectural art.
  \item A knowledge of real estate development.
  \item An insight into Washington's life.
  \item An overall review of American history.
\end{choices}
\question{23. How much should a couple with their pre-school kid pay for the Spy Museum tickets?}
\begin{choices}(2)
  \item \$33.90
  \item \$37.90
  \item \$45.90
  \item \$48.85
\end{choices}
\passagelabel{B}
\par On Dec. 26, Gelinne, 60, was gazing out his back windows of his home in Maryland, at frozen Beards Creek when he spotted a small aircraft a few hundred yards away, losing altitude.\par
\par As the plane disappeared behind the trees, Gelinne, a former Navy commander, realized it was going to land in the creek (小海湾). He flashed on a moment from more than 20 years earlier: Sept. 11,2001, he was at work when terrorists crashed a jumbo jet into the building. He fled the chaos but has always wondered if he could have stayed inside and helped.\par
\par On this day, he didn't hesitate. Racing down to the waterfront, he saw the plane had slid to a stop on the broad, frozen creek, far from the shore. It was now sinking. The pilot was standing on the wing. Gelinne knew that the icy water could kill the pilot.\par
\par Gelinne tested the ice with his foot and decided not to take any chances walking on it. So, he rushed back and swiftly pulled out a kayak(橡皮船) from his back deck.\par
\par He set off. His initial attempts left him circling on the spot. But with some muscle, he made progress. It was physically-demanding work. When Gelinne reached the plane, it had sunk; only its tail was visible. The pilot was standing on a tail wing, submerged up to his chest. Gelinne edged his kayak off the ice and into the water, paddling toward the pilot. To keep the pilot calm, he joked, "Just hang on to the boat as if you were hugging your wife." The pilot grabbed the kayak's bow. Gelinne knew he had to get the pilot out of the water immediately. But all their desperate attempts ended in failure.\par
\par By now a police officer had arrived and radioed for help. A boat from the Maryland Department of Natural Resources appeared, breaking through ice as it arrived. It picked up the pilot, Steve Couchman,71, and rushed him to safety; he would be treated for minor injuries\par
\par Rescue had arrived just in time, Gelinne recalled. "I'm 60 years old. There was no way I could get him to shore." Still, he was satisfied he'd gone the right way that day.\par
\question{24. What happened to the small aircraft?}
\begin{choices}(2)
  \item It ran out of fuel.
  \item It sank into the icy water.
  \item It flew off the course
  \item It was attacked by terrorists.
\end{choices}
\question{25. Why did Gelinne set out to help without hesitation?}
\begin{choices}(2)
  \item To make up for regret.
  \item To carry out his mission.
  \item To earn others' respect.
  \item To demonstrate his power.
\end{choices}
\question{26. Which word best describes the situation the pilot was in?}
\begin{choices}(2)
  \item Favorable.
  \item Life-changing.
  \item One-sided.
  \item Critical.
\end{choices}
\question{27. What is the best title for the story?}
\begin{choices}(2)
  \item Dragging the Pilot
  \item Braving the Elements
  \item Seizing the Chance
  \item Fighting the Currents
\end{choices}
\passagelabel{C}
\par Glass is referred to as a material which can infinitely be recycled without it impacting its quality, purity or durability. Recycled glass can be crushed into glass pieces, which can be melted down and used to produce more glass. Glass used for packaging has a high recycling rate compared to other packaging materials. In Europe, the average glass recycling rate is 76\%, compared to 41\% for plastic packaging and 31\% for wooden packaging.\par
\par When glass is left in the natural environment, it is less likely to cause pollution than plastic. Unlike plastics, which break down into microplastics that can get into soils and water, glass is non-toxic. "Glass is mainly made of silica, which is a natural substance," says Franziska Trautmann, the co-founder of Glass Half Full, a New Orleans-based company that recycles glass into sand that can be used for coastal restoration and disaster relief. Silica, also known as silica dioxide, makes up 59\% of the Earth's crust. Since it is a natural compound, there is no concern about environmental degradation.\par
\par However, glass requires higher temperatures than plastic and aluminum to melt and form, says Alice Brock, a PhD researcher at University of Southampton in the UK. Raw materials for making virgin glass also release greenhouse gases during the melting process, adding to its environmental footprint. According to the International Energy Agency, the container and flat-glass industries emit over 60 megatons of CO2 per year. A key problem with glass recycling is that it involves the remelting process, which is the most energy intensive part of glass production. It accounts for 75\% of the energy consumption during production.\par
\par Even though glass containers can be reused an average of 12-20 times, glass is often treated as single-use. Single-use glass disposed of at landfills can take up to one million years to decompose. So the next time you want to drop a glass bottle, perhaps consider reusing it first. Glass is a resilient, long-lasting material that is not made to be thrown away after only being used once.\par
\question{28. What does the passage focus on?}
\begin{choices}(2)
  \item Glass production
  \item Glass restoration.
  \item Glass prospect
  \item Glass recycling.
\end{choices}
\question{29. What does the underlined word in the second paragraph probably mean?}
\begin{choices}(2)
  \item Damage.
  \item Conservation.
  \item Assessment.
  \item Transformation.
\end{choices}
\question{30. What can we learn from paragraph 3?}
\begin{choices}(2)
  \item Glass melting process is complicated.
  \item Glass reduces environmental footprint.
  \item Glass items are meant for permanent use.
  \item Glass production consumes much energy.
\end{choices}
\question{31. Why does the author advocate reuse of glass containers?}
\begin{choices}(2)
  \item To confirm their durability.
  \item To suggest being economical.
  \item To promote the green idea.
  \item To propose a new energy plan
\end{choices}
\passagelabel{D}
\par Exhausted by the rising cost of living in the United States and non-stop advertisements, some young adults on TikTok are pushing back.\par
\par "When every moment of your life feels like you are being sold something and the price of that item keeps going up, people will burn out on spending money," said Kara Perez, an influencer and financial educator.\par
\par Social media has long showcased picture — perfect homes, luxurious closets and various beauty products. But a new trend, known as "underconsumption core", is moving in the opposite direction — encouraging more economical lifestyles and prioritizing quality over quantity. For example, many people on social media proudly display their well-worn shoes that have been with them since middle school, along with their commitment to using the last drop of their makeup. Others opt for repurposing old mason jars instead of buying new food storage containers.\par
\par For a long time, young people in the West have commonly made use of "Buy Now, Pay Later" plans, which have intensified consumption. They have developed a behavior of spending down to their last pound on a fashion item. It is an addiction tied to the pressure "to show who we are through possessions". However, years of inflation(通货膨胀) have forced many to conclude that they cannot keep up with the spending habits of those on their social media feeds.\par
\par "The social media trend of underconsumption' is a way for Gen Z to make the most of their money and be environmentally friendly at the same time," said Ashley Ross, head of Bank of America.\par
\par Low-consumption brick-and-mortar (实体的) initiatives attract a broader range of people across different profiles and generations. Anjali Zielinski,42, joined a "Mending 101" workshop in Georgetown, DC, hoping to acquire new skills. She brought her 7-year-old daughter, Mina, along with her. In addition to providing an outlet for Mina's creativity, she hopes the craft will teach her daughter the "value of our possessions and the work that goes into them."\par
\par But some are still doubtful. It remains to be seen whether "underconsumption core" is simply a mood or a lasting movement. "The lifespan of this trend depends, in many ways, on how long the economy continues to slow and incomes remains below price gains," Brett House, an economics professor at Columbia Business School, told CNBC.\par
\question{32. What is the cause of some Americans' cutting back on spending?}
\begin{choices}(2)
  \item The commercial activities.
  \item Their social media feeds.
  \item Their material satisfaction.
  \item The economic downturn.
\end{choices}
\question{33. What is the philosophy of "underconsumption core"?}
\begin{choices}(2)
  \item Spending money as you like
  \item Doing more with less.
  \item Keeping up with the Joneses.
  \item Buy Now, Pay Later.
\end{choices}
\question{34. How does the author tell us the trend's popularity among people?}
\begin{choices}(2)
  \item By giving examples.
  \item By making comparisons.
  \item By listing numbers.
  \item By citing research data.
\end{choices}
\question{35. What is Brett House's attitude to the "underconsumption core" trend?}
\begin{choices}(2)
  \item Doubtful.
  \item Reserved.
  \item Objective.
  \item Disapproving.
\end{choices}
\sectitle{第二节（共5小题；每小题2.5分，满分12.5分）}
\par 阅读下面短文，从短文后的选项中选出能填入空白处的最佳选项。选项中有两项为多余选项。\par
\par About 90 percent of the world's population doesn't need to think about which hand to use when picking up a pair of scissors. \blanknum{36} Let's take a look at some of the myths, mysteries and history of our left-handed friends.\par
\par \blanknum{37}. The belief that left-handed people were witches was widespread in medieval Europe, and they could even be killed. In fact, the word "sinister," meaning harmful or evil, comes from a Latin word for "on the left side." Even in the 1960s, American children born left-handed were forced to switch hands when writing.\par
\par There are also some interesting mysteries concerning right- or left-handed people. The human brain is "cross-wired" with the left brain controlling the right side of the body and vice versa(反之亦然). But one University of Oxford study suggests that, in the brains of left-handed people, the two parts are better connected. \blanknum{38}.\par
\par Unfortunately, there are some mental health issues that can especially affect left-handed people. According to a 2013 Yale University study, lefties seem more likely to suffer from schizophrenia(精神分裂症). \blanknum{39}. But other studies have linked the left-handed brain with decidedly creative thinking. Though uncertain, yet with artists such as Leonardo da Vinci, there might be some connection.\par
\par One area where lefties may have the upper hand is in sports that require quick reaction times. \blanknum{40}. Ding Ning, for example, is a left-handed table tennis champion with three Olympic gold medals. But if there is an advantage, it's probably not a result of some special "lefty magic." Recent studies show that there seems to be no definite connection being left-handed and athletic capability.\par
\par In any case, lefties, hold your southpaw high and show the world what you're made of!\par
\begin{choices}(1)
  \item Altogether, it's not a shame to be left-handed
  \item Being left-handed hasn't always been accepted
  \item This includes sports like baseball and table tennis
  \item It is this disease that has led to their being left-handed
  \item But it's not the same story for the 10 percent who are left-handed
  \item Also, studies say that lefties may have an advantage in language learning
  \item This serious disorder can cause people to hear voices or see things that are not real
\end{choices}
\parttitle{第三部分 语言知识运用（共两节，满分 30 分）}
\sectitle{第一节（共15小题；每小题1分，满分15分）}
\par 阅读下面短文，从短文后各题所给的 A、B、C 和 D 四个选项中，选出可以填入空白处的最佳选项。\par
\par In Summersville, where the sun always shone, young Tommy, a \blanknum{41} enthusiast, set out to adventure under the sun. His mission was simple: to \blanknum{42} every hidden gem and secret spot around town.\par
\par The town was boiling hot, but Tommy's spirit was as \blanknum{43} as a cucumber. He studied an old map and decided to \blanknum{44} an old lighthouse. It had been \blanknum{45} for years. And to add to its mystery, the lighthouse was said to be haunted (闹鬼). Tommy hesitated, but in the heat of moment, he decided to \blanknum{46}.\par
\par Inside, he discovered a hidden time capsule. The walls were decorated with paintings that \blanknum{47} Summersville's amazing history. Each painting was a story of days filled with \blanknum{48} and nights with magic. Then and there, Tommy realized his summer adventure was about \blanknum{49} the soul of his town.\par
\par Tommy committed his summer to \blanknum{50} the lighthouse. He made up his mind to turn the lighthouse into a place of beauty for all to \blanknum{51}. His efforts brought the community together, transforming the \blanknum{52} lighthouse into a beautiful landmark.\par
\par Autumn approaching, the lighthouse stood tall and proud. Tommy's \blanknum{53} had become a gift to Summersville. And as the sun set on the last of the summer days, Tommy sat by the shore, his heart \blanknum{54}. His summer had proved to be more \blanknum{55} than he could have ever imagined.\par
\cz{41}{sports}{summer}{travel}{charity}
\cz{42}{publicize}{introduce}{advertise}{explore}
\cz{43}{cool}{long}{rough.}{fresh}
\cz{44}{run into}{set up}{check out}{take over}
\cz{45}{abandoned}{employed}{established}{maintained}
\begin{choices}(3)
  \item quit 3. investigate
  \item wait
  \item withdraw
\end{choices}
\cz{47}{denied}{dominated}{illustrated}{emphasized}
\cz{48}{sorrow}{relief}{hurry}{joy}
\cz{49}{pursuing}{appreciating}{hiding}{uncovering}
\cz{50}{constructing}{restoring}{visiting}{guarding}
\cz{51}{enjoy}{record}{expand}{shift}
\cz{52}{ordinary-looking}{splendid}{once-forgotten}{unique}
\cz{53}{adventure}{painting}{lighthouse}{map}
\cz{54}{anxious}{broken}{full}{strong}
\cz{55}{demanding}{romantic}{wasteful}{exciting}
\sectitle{第二节（共10小题；每小题1.5分，满分15分）}
\par 阅读下面短文，在空白处填入 1 个适当的单词或括号内单词的正确形式，\par
\par Honesty is the best policy. But in many situations, people do not want the truth. Whether it's a fact of life \blanknum{56} a criticism of their behavior, everyone has truths they want to ignore. These truths \blanknum{57} (call) truth bombs. This metaphor shows truth can cause damage if not \blanknum{58} (handle) well, so caution should be exercised when pointing out difficult truths\par
\par Sometimes we drop truth bombs out of anger. We may rationalize(合理解释) our words by saying that we need to express our \blanknum{59} (emotion) or to help change their ways. But often we use the truth as \blanknum{60} weapon to make the other person feel bad and to repay them for what they did\par
\par Truth bombs rarely lead to \blanknum{61}(production) discussions because they trigger(触发)a defensive reaction. Over time, truth bombs can increase conflict and weaken trust in relationships. Also, truth bombs often fail \blanknum{62}(represent) the whole truth. Anger can make it difficult to see the big picture.\par
\par Sometimes people do need to hear hard truths. Pretending that problems do not exist can weaken trust just as much as \blanknum{63} (drop) truth bombs can. And failing to address problems in relationships can lead to resentment(怨恨), which may eventually boil over into conflict. By presenting truth and addressing problems \blanknum{64} a constructive way, you can strengthen relationships and avoid the emotional damage that comes from using \blanknum{65} (over) blunt(直言的) words.\par
\parttitle{第四部分 写作（共两节，满分 40 分）}
\sectitle{第一节（满分15分）}
\par 假如你是李华，你的英国好友 Eric 酷爱音乐创作并刚跟你分享了他的新歌"Dancer"，现请你给它写一\par
\par 邮件以表欣赏，内容要点包括：\par
\par (1) 所听及感受;\par
\par (2) 期待更多作品。 注意：\par
\par (1) 写作词数应为 80 左右;\par
\par (2) 请按如下格式在答题纸的相应位置作答。\par
\par\vspace{0.4em}
\begin{center}\begin{tabular}{{|l|}}
\hline
Dear Eric, \\
Yours, Li Hua \\
\hline
\end{tabular}\end{center}
\vspace{0.4em}
\sectitle{第二节（满分25分）}
\par 阅读下面材料，根据其内容和所给段落开头语续写两段，使之构成一篇完整的短文，\par
\par Anika pulled her wheeled suitcase behind her, following her parents through the airport. It would be her first time on an airplane. She was very excited and somewhat nervous as well. She hugged her favorite doll tightly. That made her feel better.\par
\par She boarded the plane and walked down the narrow aisle 通道 to find her seat with her parents. A flight attendant smiled at her and asked where they were going. "Walt Disney World," Anika responded with a smile.\par
\par "You will have a wonderful time." The attendant pointed to her name tag. "I'm Leanne. You let me know if you and your pretty doll friend need anything."\par
\par "Thanks. I'm Anika, and this is Pansy," said Anika, holding the doll up to show her. Leanne stroked Pansy and left to attend to other passengers.\par
\par The flight was going to be quite a long time. Anika read a book with Pansy and did some puzzles her mother had packed for the trip. Soon she was getting sort of tired. Anika's mom suggested her taking a nap.\par
\par They asked Leanne for a blanket. "It should keep you warm and cozy," Leanne said. "Night-night!" She winked at the young girl. Anika tucked Pansy under the blanket with her. Soon, she fell asleep against her father's shoulder.\par
\par The plane touched down in Florida and Anika was woken up. Her parents collected the luggage and headed to the exit. Anika followed, not quite awake yet. At the door, Leanne waved them goodbye.\par
\par All passengers had left and Leanne began to gather blankets from the seats. Under one blanket, she was surprised to find a doll. She picked it up and recognized it was Pansy. She thought about catching up with the family, but soon she dismissed the idea, for it was too late. She just hoped to connect with the little girl again in some way. Then she put away Pansy for safe keeping. 注意:\par
\par (1) 续写词数应为 150 左右;\par
\par (2) 请按如下格式在答题卡的相应位置作答。\par
\par\vspace{0.4em}
\begin{center}\begin{tabular}{{|l|}}
\hline
Paragraph 1: When the family arrived at the hotel, Anika suddenly remembered Pansy.\_ \\
Paragraph 2: They boarded the plane and Anika was surprised to see Leanne again.\_ \\
\hline
\end{tabular}\end{center}
\vspace{0.4em}
\end{document}
